% Nguồn: src/react_agent/security_v3/{intent_policy_v1,value_labels_v1,param_pinning_v1,brokerguard_v3_v1}.py;
% decision log 2026-11-16 12:30 (C1), 15:30 (runtime, bỏ C1a), 16:20 (vi mô), 2026-11-17 19:25 (đóng băng).
% Số: facts.md mục "Bổ sung cho bài v3" (ngưỡng 8 ký tự, 9 chữ số; vi mô) và mục Test-4 (0,26 ms). Bản 15 trang bỏ
% đoạn về biến thể C1a (2 -> 5 quỹ đạo gây hại trên Dev-val), xem bản dài ở commit 6814b3a.
% Không trích văn bản tấn công. Ví dụ từ nối phủ định là của bộ trích, không phải của dữ liệu.
\section{Phương pháp \ours}\label{sec:method}

Thiết kế của \ours{} xuất phát từ quan điểm rằng mô hình không phải là ranh giới an ninh. Một hành động chỉ có hiệu lực
khi broker thực thi nó, nên broker đóng vai trò reference monitor tại các sink, và nguồn quyền duy nhất là $u$. Lời gọi
bị từ chối được báo lại để tác tử làm tiếp phần việc hợp lệ. Hình~\ref{fig:arch} mô tả kiến trúc. Các luật nền có từ
giai đoạn đầu gồm \rDest{} (Mục~\ref{sec:intro}), \rFlow{} (giá trị từ nguồn gắn nhãn mật chỉ được rời tác tử khi $u$
chứa nguyên văn giá trị đó), \rRedact{} cho câu trả lời cuối và \rSig{} (chặn sink sau một quan sát bị bộ dò cụm từ
gắn cờ, trừ lời gọi khớp một văn phạm chặt trên $u$). Lời nhắc kiểu sandwich nhắc lại $u$ sau mỗi quan sát. \ours{} thay
\rDest{} bằng \pOne{} (\cOne), thêm \pTwo{} (\cTwo) và \pThree{} (\cThree). Cấu hình giữ \rDest{} và bỏ cả ba thành
phần được gọi là \strSand, khi bỏ thêm lời nhắc thì gọi là \strRules.

% Hình 1, TikZ (Overleaf tự vẽ, không render trên máy). Tọa độ tính bằng cm, tổng rộng 11,85 cm < \textwidth 12,2 cm.
% Ô kiểm 1,7 x 0,9 cm, chữ 6,4 pt: tối đa khoảng 14 ký tự mỗi dòng, ba dòng mỗi ô (giữ nhãn ngắn khi sửa).
% Màu theo thực thể (giống Hình 2, 3): cam viền đậm = thành phần mới C1--C3, xanh = giữ từ BG-Base, xám nét đứt = phụ
% thuộc mô hình. In đen trắng: phân biệt bằng độ dày viền và nét đứt.
% Nhãn tiếng Anh cho bản nộp:
%   User / $u$ (trusted) | LLM agent / (ReAct) | Tool Broker (host side)
%   C1: extract $P(u)$ from $u$ / (destination, polarity, pinned values) | C2: label each value in $o_t$ / (protected set $C_t$)
%   before every sink call $a_t$: executed only if all four checks pass
%   R-Signal / flagged, so deny | C1 / positive grant | C3 / pinned values kept | C2 + R-Flow / no value of $C_t$ leaves
%   R-Redact + C2 / redact values in $y$ | Sandwich / repeat $u$ | Report denial, / agent continues
%   Needs confirmation / = denial | Read tools / (untrusted) | Sinks / e-mail, webhook | $y$ redacted
%   new component | kept from BG-Base | model-dependent
\begin{figure}[!htbp]
\centering
\begin{tikzpicture}[
  font=\fontsize{6.4}{7.6}\selectfont,
  box/.style={rounded corners=1pt, align=center, inner sep=1pt},
  newc/.style={box, draw=bgorange, line width=0.9pt, fill=bgorange!12},
  kept/.style={box, draw=bgblue, line width=0.45pt, fill=bgblue!10},
  soft/.style={box, draw=black!55, line width=0.5pt, fill=black!4, dash pattern=on 2pt off 1.3pt},
  neutral/.style={box, draw=black!60, line width=0.5pt, fill=black!5},
  arr/.style={-{Stealth[length=3pt]}, line width=0.5pt, draw=black!65},
  lab/.style={font=\fontsize{6.4}{7}\selectfont},
  note/.style={font=\fontsize{5.6}{6.6}\selectfont\itshape, text=black!65, anchor=west},
  gate/.style={minimum width=1.7cm, minimum height=0.9cm, text width=1.6cm},
  swatch/.style={minimum width=0.34cm, minimum height=0.24cm},
]
% cột trái: người dùng và tác tử
\node[neutral, minimum width=1.5cm, minimum height=0.8cm, text width=1.4cm] at (0.75,4.85)
  {Người dùng\\ $u$ (đáng tin)};
\node[neutral, minimum width=1.5cm, minimum height=1.3cm, text width=1.4cm,
  font=\fontsize{6.6}{7.6}\selectfont\bfseries] at (0.75,2.35) {Tác tử LLM\\ (ReAct)};
\draw[arr] (0.75,4.45) -- (0.75,3.0);

% khung broker
\draw[rounded corners=2pt, line width=0.6pt] (2.0,0.5) rectangle (9.3,5.65);
\node[font=\fontsize{7}{8}\selectfont\bfseries] at (5.65,5.43) {Tool Broker (phía host)};

% hàng 1: trạng thái do C1 và C2 dựng
\node[newc, minimum width=3.45cm, minimum height=0.66cm, text width=3.35cm] at (3.825,4.8)
  {\textsf{C1}: trích $P(u)$ từ $u$\\ (đích, cực tính, giá trị ghim)};
\node[newc, minimum width=3.45cm, minimum height=0.66cm, text width=3.35cm] at (7.475,4.8)
  {\textsf{C2}: gắn nhãn từng giá trị trong $o_t$\\ (tập bảo vệ $C_t$)};
\draw[arr] (1.5,4.8) -- (2.1,4.8);

% hàng 2: bốn kiểm tra trước mỗi lời gọi sink
\node[note] at (2.1,4.2) {trước mỗi lời gọi sink $a_t$: chỉ thực thi khi cả bốn kiểm tra đều qua};
\node[kept, gate] at (2.95,3.5) {\textsf{R-Signal}\\ quan sát bị gắn cờ thì chặn};
\node[newc, gate] at (4.75,3.5) {\textsf{C1}\\ đích có quyền dương};
\node[newc, gate] at (6.55,3.5) {\textsf{C3}\\ giá trị ghim còn nguyên};
\node[newc, gate] at (8.35,3.5) {\textsf{C2} + \textsf{R-Flow}\\ không lộ giá trị của $C_t$};

% hàng 3: câu trả lời cuối, lời nhắc, phản hồi, xác nhận
\node[kept, gate] at (2.95,2.1) {\textsf{R-Redact} + \textsf{C2}\\ che giá trị trong $y$};
\node[soft, gate] at (4.75,2.1) {Sandwich\\ nhắc lại $u$};
\node[kept, gate] at (6.55,2.1) {Báo từ chối,\\ tác tử làm tiếp};
\node[newc, gate] at (8.35,2.1) {Cần xác nhận\\ = từ chối};

% chú giải
\node[newc, swatch] at (2.3,0.88) {};
\node[note] at (2.52,0.88) {thành phần mới};
\node[kept, swatch] at (4.3,0.88) {};
\node[note] at (4.52,0.88) {giữ từ \base};
\node[soft, swatch] at (6.3,0.88) {};
\node[note] at (6.52,0.88) {phụ thuộc mô hình};

% tác tử <-> broker
\draw[arr] (1.5,2.75) -- node[lab, above=-1pt] {$a_t$} (2.0,2.75);
\draw[arr] (2.0,2.35) -- node[lab, above=-1pt] {$o_t$} (1.5,2.35);
\draw[arr] (1.5,1.95) -- node[lab, above=-1pt] {$y$} (2.0,1.95);

% cột công cụ
\node[neutral, minimum width=2.1cm, minimum height=0.8cm, text width=2.0cm] at (10.8,3.95)
  {Công cụ đọc\\ (không đáng tin)};
\node[neutral, minimum width=2.1cm, minimum height=0.8cm, text width=2.0cm] at (10.8,2.45)
  {Sink\\ e-mail, webhook};
\draw[arr] (9.3,4.15) -- node[lab, above=-1pt] {$a_t$} (9.75,4.15);
\draw[arr] (9.75,3.75) -- node[lab, above=-1pt] {$o_t$} (9.3,3.75);
\draw[arr] (9.3,2.45) -- node[lab, above=-1pt] {$a_t$} (9.75,2.45);

% câu trả lời đã che rời broker
\draw[arr] (9.3,1.0) -- (9.95,1.0);
\node[anchor=west, text=black!65] at (9.98,1.0) {$y$ đã che};
\end{tikzpicture}
\caption{Kiến trúc của \ours. Mọi lời gọi $a_t$ và câu trả lời cuối $y$ đều đi qua Tool Broker. \cOne{} trích chính
sách $P(u)$ một lần từ $u$, còn \cTwo{} cập nhật tập giá trị bảo vệ $C_t$ sau mỗi quan sát. Ô viền cam đậm là thành
phần mới, ô viền xanh giữ từ \base, ô nét đứt phụ thuộc vào mô hình. Mọi kiểm tra trong khung đều tất
định.}\label{fig:arch}
\end{figure}


\paragraph{Chuẩn hóa.} Hàm $\norm(\cdot)$ áp Unicode NFKC, bỏ ký tự định dạng và casefold. Hàm $\kappa(\cdot)$ áp thêm
việc bỏ dấu tiếng Việt và chỉ giữ chữ và số. Ký hiệu $x \sub z$ nghĩa là $\kappa(x)$ là một chuỗi con của $\kappa(z)$,
và $V(u)$ là tập giá trị xuất hiện trong $u$ theo nghĩa này.

\paragraph{\POne{} (\cOne).} Bộ trích chỉ đọc $u$ và trả về $P(u) = \{g_i\}$. Mỗi quyền $g = (\mathrm{tool}, \dest, \pi,
\mathrm{pins})$ gồm công cụ, đích, cực tính $\pi \in \{+,-\}$ và các giá trị được ghim. Đích là địa chỉ e-mail hoặc
endpoint, so khớp nguyên token sau $\norm$, nên một phần của đích hay một từ thường trong $u$ không còn được coi là
đích. Phủ định đi với động từ (``đừng gửi'', ``do not send'') phủ các đích theo sau trong cùng mệnh đề, từ loại trừ
(``trừ'', ``except'') phủ tới dấu phẩy kế tiếp, còn từ tương phản (``nhưng'', ``but'') kết thúc phạm vi. Khi $u$ không
dấu, một từ như ``dung'' chỉ được coi là phủ định nếu theo sau là một động từ gửi. Mệnh đề cấp quyền phải có động từ gửi
không bị phủ định và hợp với loại đích. Một đích chỉ được nhắc tới, nằm trong trích dẫn hay đứng gần một phủ định không
rõ phạm vi cho kết quả \emph{cần xác nhận}. Lời gọi sink $a$ thỏa \pOne{} khi
\begin{equation}\label{eq:c1}
\begin{aligned}
&\exists\, g \in P(u)\colon\; \pi(g) = {+} \,\land\, \mathrm{tool}(g) = \mathrm{tool}(a) \,\land\, \dest(g) \doteq \dest(a),\\
&\nexists\, g' \in P(u)\colon\; \pi(g') = {-} \,\land\, \dest(g') \doteq \dest(a),
\end{aligned}
\end{equation}
với $\doteq$ là bằng nhau nguyên token sau $\norm$. Chẳng hạn, với yêu cầu ``gửi báo cáo cho \code{ketoan@example.com},
không gửi cho \code{doitac@example.org}'', lời gọi tới địa chỉ thứ hai bị từ chối dù \rDest{} cho qua.

\paragraph{\PTwo{} (\cTwo).} \PTwo{} gán nhãn cho từng giá trị trong mọi đầu ra công cụ theo kiểu của nó, độc lập với
nhãn nguồn và độ tin cậy. Ba kiểu được nhận ra là định danh trộn chữ và số dài ít nhất 8 ký tự, số có ít nhất 9 chữ số
sau khi bỏ khoảng trắng, dấu chấm và gạch, và giá trị đứng sau tên trường nhạy cảm như mật khẩu, OTP, token hay số tài
khoản. Gọi $C_t$ là tập giá trị cần bảo vệ sau bước $t$, gồm các giá trị được \rFlow{} bảo vệ và các giá trị được gắn
nhãn theo kiểu. Broker từ chối lời gọi sink $a$ nếu
\begin{equation}\label{eq:c2}
\exists\, v \in C_t \setminus V(u)\colon\; v \sub \data(a) .
\end{equation}
Vì so khớp dùng $\kappa$, điều kiện vẫn đúng khi giá trị bị thêm hay bớt khoảng trắng, dấu chấm, gạch, đổi hoa thường
hoặc bỏ dấu. Trong câu trả lời cuối, mọi $v \in C_t \setminus V(u)$ có $v \sub y$ được thay bằng \code{[REDACTED]}.

\paragraph{\PThree{} (\cThree).} Sửa tham số giữ nguyên đích được cấp quyền và đổi một giá trị như số tiền, trạng thái,
đơn vị hay mã. Thành phần này ghim các literal trong mệnh đề của $u$ đã cấp quyền cho lời gọi, gồm số đứng sau động từ
gửi và không phải toán hạng của một phép tính, đơn vị nêu kèm số, giá trị boolean có cực tính và định danh do người dùng
viết. Lời gọi bị từ chối khi thiếu một literal đã ghim trong $\data(a)$, so ở dạng $\kappa$, hoặc mang boolean ngược với
boolean đã ghim.

\paragraph{Cổng và chi phí.} Một lời gọi sink chỉ được thực thi khi qua cả bốn kiểm tra là \rSig, \pOne, \pThree{} và
kiểm luồng gồm \rFlow{} cùng \pTwo. Kết quả \emph{cần xác nhận} được thi hành như từ chối và đếm riêng, còn lời gọi đọc
luôn được cho qua. Cổng không gọi thêm mô hình nào. Trên tập phát triển, trung vị thời gian kiểm một lời gọi sink là
315~$\mu$s (\strSand: 107~$\mu$s).
